\section{Objective Reality}

\begin{quotex}
Our age is an age in which men rise painfully from effects to causes, in which they even concern themselves only with effects, in which they say it is useless to concern oneself with causes, and in which they do not even know what a cause is.
\flright{\textsc{Joseph de Maistre}}

The crown of the Emperor signifies the renunciation of freedom of intellectual movement, just as his arms and legs signify his renunciation of freedom of action and movement. He is deprived of the three so-called ``natural'' liberties of the human being---those of opinion, word and movement. Authority demands this.
\flright{\textsc{Valentin Tomberg}}
\end{quotex}


Every possible world requires five conditions:

\begin{itemize}


\item Form and matter

\item Space and Time

\item Consciousness (or Life)
\end{itemize}


A world without consciousness is inconceivable; try it if you can. Space and Time are necessary for conscious experience. Although they may be experienced differently in different worlds, there needs to be some sense of synchrony and diachrony.

Anything that exists in the world requires matter and form, or essence. These two factors are related to our knowledge of the external world and the intelligibility of the world.

\begin{itemize}


\item The objective reality of the world is imposed on us through the senses

\item The ``crown of thorns'' give conscience to objective thought

\item Thought, in its pure state, is the organ of truth, not illusion

\item However, the five dark currents of the will take hold of the thinking centre, thereby distorting it
\end{itemize}


\paragraph{The External World}


\begin{quotationx}
A wound is a door through which the objective exterior world intrudes into the interior of the closed system of the subjective interior world. Speaking biologically, it is a breach in the walls of the fortress of the organism by which forces from outside the organism penetrate into its interior.

The eyes are open wounds which are so sensitive that they react to every nuance of light and colour. And it is the same with the other sense organs. They are wounds, i.e. it is they which impose on us the objective reality of the outer world. Where I would like to see beautiful flowers, my eyes make me see a pile of dung. I am forced to see what the objective world shows me by way of my eyes. It is like a nail from outside nailing my will.

\emph{The senses---given that they are sound and functioning normally---are wounds through which the objective world, without regard to our will, imposes itself on us.}
\flright{\textsc{Valentin Tomberg}}
\end{quotationx}


Hence, we create a representation of the external world in consciousness, seemingly against our will. Moral action requires a Not-I, in which there are objects in opposition to the I, and which pose a challenge to our projects. A weak will may deny aspects of the external world, and a strong will accepts the challenge.

\paragraph{The Intelligibility of the World}


\begin{quotationx}
The head bears the ``crown of thorns'', which is borne, in principle, by every person capable of objective thought --- the ``crown of thorns'' being given to the human being since the beginning of human history. It is that subtle organ which is designated as the crown centre. This crown centre is a ``natural gift'', as it were, to each human being and every normal person possesses it. The ``thorns'' of the crown centre function as the ``nails'' of objectivity, which give conscience to thought. It is thanks to them that thought has not become wholly emancipated and as arbitrary, for example, as the imagination is.

\emph{Thought as such is, in spite of all, the organ of truth, not of illusion.}
\flright{\textsc{Valentin Tomberg}}

A process of mind is thinking only so far as it realizes the end of thinking, which is always to understand, that is, to see things within a system which renders them necessary.
\flright{\textsc{Brand Blanshard}, \emph{The Nature of Thought}}
\end{quotationx}


The senses show that things exist, but not \emph{what} exists. In other words, the senses are open to the material condition of the world and the mind perceives essences. Initially it grasps the essences of individual things; this is so common that few ever become aware of it.

Living thought then sees things in relationship to each other, both synchronically and then diachronically, in larger and larger wholes, until ultimately it grasps the Absolute. This is spiritual vision and enables one to see the future.

Contrary to popular belief, no mere accumulation of facts can lead to an understanding of the world; only objective thought can do so. Facts do not create theories. This is most emphatically illustrated in the systems created first by Newton and then by Einstein. Newton's theory, despite its effectiveness in predicting events in the external world, was replaced by Einstein's more comprehensive understanding. So if Newton's theory was NOT dictated by the facts, it must therefore be a free creation of the human spirit. The direction is one way: facts don't create the theory, but the theory creates the facts.

If thinking is the organ of truth, then how can we account for all the divergent differences in thought and opinions?

\paragraph{Evil Will}


\begin{quotex}
Imagine that the five organs of action---the limbs, including the head in its function as a limb---were to have analogous wounds, i.e. that the five currents of will of which they are an expression were to give access to an objective will which would be to personal desires what sense perceptions are to the play of fantasy.
\flright{\textsc{Valentin Tomberg}}
\end{quotex}


When functioning properly, the mind discerns the good and guides the will. Otherwise, the will controls the organ of thinking, using it merely as its instrument. Thinking, in this case, will then justify what the will desires; it sinks into subjectivity and fantasy, both opposed to its God-given function. Irrationality---that is, the denial of the Logos and cosmic order---and outbursts of anger are its bad fruits.

Those are the marks of an evil will.

\paragraph{Cause and Effect}


People tend to work backwards, from effects back to causes. As Hume demonstrated, however, this concept of cause does not arise through reason, but follows rather from a habitual way of looking at things. This is another way of saying the facts cannot create understanding, so knowledge of effects cannot dictate the knowledge of causes.

Nevertheless, through force of habit, men fantasize about causes in order to try to understand the world. Hence, in the past, the unknown substance phlogiston was assumed to create fire and spontaneous generation to explain maggots in garbage. In our time, there is no shortage of fanciful causes to explain events. One group will claim that it is white privilege or the patriarchy; another will blame Christianity or the Jews. Books are written to discover the hidden causes behind events. They are popular because men prefer easy answers, and it gives them the illusion of intelligence.

What is much more difficult, however, is to discern the effects, especially the long-term effects of causes. The effects of a policy embarked on today will usually not be experienced for a generation or more. It takes a rare intelligence to anticipate the future undesired consequences of what is done today.

\textbf{Joseph de Maistre} was not interested in the logic of facts, since, as he asserted, the Law of Gravity does not cause the rock to fall. Instead, he was concerned with the moral logic, or the moral cause of events. Science deals with facts yet can never answer the question, ``Why?'' Ultimately, therefore, it comes down to the moral will, which is the real cause. And an evil will results in evil effects.

A bullet or a gun is not the cause of murder, but rather the person who pulled the trigger.

\paragraph{The Spirit of Schism}


People take great pride in causing division through partisanship\footnote{\url{https://gornahoor.net/?p=40}}. In the West, for example, there are repeated examples of schisms that broke up the spiritual unity of Europe. The illusion is that there are only two choices. Ultimately, another option comes into play, if for no other reason than to end the bloodshed of the warring factions. Even though everyone knows that ``a house divided against itself cannot stand'', the will-to-greatness drives men to settle for half of the house.


\flrightit{Posted on 2019-05-12 by Cologero}

\begin{center}* * *\end{center}

\begin{footnotesize}
\begin{sffamily}
\texttt{Han Fei on 2019-05-16 at 23:28 said:}

Hello Cologero, I wish to ask you this time about ``black pills''? How does one address the creeping feeling that everything is getting worse and worse by the year, that the ability in people to perceive right from wrong is fading by the day, that all that is noble in man is being slowly replaced by spite, arbitrary vice and hatred. Is it true that the further one immerses into the world, the more of the evil, insufferable aspect of the reality is disclosed to him along the way -- is this what the early Church fathers meant by their seemingly inexplicable denunciation to scientific means of enlightenment and pursuit of worldly goals? Whenever I encounter society around me, I do not perceive human beings as evil by nature and as such deserving to be hated, even as they utter the most soul crushing words, but they are misguided, driven into ruin by their leaders, engorged in madness by popular culture. I have to keep on reminding myself of this, otherwise I might fall into the same trap.

\hfill

\texttt{Cologero on 2019-05-17 at 17:00 said:}

Han Fei, good points. I was going to respond, but it is turning into a blog post, hopefully this weekend.

\hfill

\texttt{Fallen Adam on 2019-05-17 at 19:17 said:}

Can consciousness be described as being the Form of a being, when it is brought into actuality on the plane of creation, as in, that it's the fulfilling of a beings form and final cause in creation, when it lives according to it's nature? Using the terminology Mouravieff uses, is conscience the neutralizing principle between the active principle of the form, and the passive one of matter? Basically, i'm asking if consciousness is the flowing of the essence of a being in time.

\hfill

\texttt{Han Fei on 2019-05-17 at 21:28 said:}

From the way I worded things, it would appear that some of the articles already posted here address this question. But I want to make it clear that I do not sense the world to be any way ``evil'' in itself. What is expressed here is a feeling of exasperation at the human element of the world, that is society, the creeping feeling of the falsity of men's motives, and of the painful realization of the lack of sincerity in societies and institutions which in principle should be the grounds of them.

This conundrum would be best served by an example. In one city the authorities want to build a large church. The cathedral will grandiose and beautiful in stark contrast to the commercial and utilitarian architecture of this day. But it will be erected in the last park the city has, cutting its accessible area by around half. Naturally this draws huge protests from the city denizens for whom this park presents the last area of greenery in downtown. A rather ironic picture emerges -- to save a grove of trees or build (or in the case I speak of, rebuild) a temple of God.

The fight that the citizens to put up against the, perfectly lawful I must add construction, isn't really about a couple of trees, but because we perfectly know that a church provides limited interest to the majority of the people who are not believers, where as a park represents a universally recognized practical value. From the perspective of one side, it is part of a state undertaken project to revitalize its traditional identity by sanctifying the historical center of the town and from the perspective of another, nothing but a cynical ploy to push state ideology of corruptocrats using ``insane people who believe in fairy tales'' (a statement repeated ad nauseam on social networks discussing this event) spending billions of dollars doing so. Both sides are right in their own way, and the reasons behind this interminable opposition is precisely the cause of my distress at the state of the world.

Of course this might all sound like a trivial reason to get worked up about. But let's not forget that a variation of this same tale left heads flying in other times and places.

\hfill

\texttt{Fallen Adam on 2019-05-19 at 06:15 said:}

I've thought a little bit more about this question, and i can now express better what i meant. According to Thomism, every created thing must be created by an rational agent, being made towards a certain end, which is that created being final cause. If that being upon being created, already had his final cause ``accomplished'' so to speak, then is existence in creation wouldn't really be different from his essence or archetype. But if that created being had no causal or creative power of his own, only being acted on by his creator, then he wouldn't really be a distinct being, and there would be no reason to call him a whole, as any action done by it would need to be moved by his creator. Thereby, besides the form and matter which compose a being, there must be another principle to account for their living and acting on the realm of creation, which is consciousness and life. I don't know if there is any mistake in here, or if i have forgotten anything.

In regards to the question of knowledge, our intellect can grasp the form of things or their archetypes, there being an union between the subject and object, actualizing the potentialities of the intellect, when it comes to higher states of being, but when it comes to conscience, we can't really ``have'' the conscience of an angel, in the sense that we can't stop being a man, even if we actualize those potentialities of the intellect. So when it comes to conscience, it seems that the way for us to understand the life or conscience of other types of beings, is trough an analogy with forces and processes that happen within us, with this idea being given validity by the fact that man contains the macrocosm in the microcosm. Can it be said then, that while knowledge of forms relates to the intellect, the understanding of conscience relates to the will? One can think here of the subject of magic, which is this control and directing of the forces of consciousness and life in creation, and which thereby is done by the will.

\end{sffamily}
\end{footnotesize}

